\subsection{Discussion}
\label{Discussion3}
\setlength{\parindent}{20pt}
\justifying

\noindent
Several design modifications aimed to make this third version of the experiment easier for participants. The conversion of BACS-1 stimuli into Latin characters and the consequent removal of the morpheme spelling task certainly succeeded in increasing word segmentation ability, as this was reflected in both higher accuracy and faster reaction times in the morpheme familiarity task. These modifications seem to have made the stimuli simpler to process. That no multilingual metrics impacted segmentation ability might indicate that the stimuli might have been learned sufficiently well, with the main brunt of cognitive processing being shifted to the tracking of other regularities. Segmentation at this level of difficulty may not require different cognitive strategies developed from multilingualism.

A higher probability of classifying stimuli as congruent arose for the male participants, and for those multilinguals who were more balanced across all of their languages. Gender differences as a main effect were only found in this last version of the experiment, and fall outside the scope of this study, and will not be discussed further. On the other hand, participants engaged in regular switching and/or a more egalitarian view of their languages might have more fluid cross-linguistic transfers, which might lead them to accept a wider range of stimuli as linguistic.

The morpheme familiarity effect at group-level saw a downward shift compared to the previous versions of the experiment. The average proportion of ``yes'' responses in the 0M were at 38\%, with the 1M condition taking its place at 53\%. Responses to the 2M condition were far above this at 68\%. This could be due to the aforementioned change in stimulus nature, and the consequent rise in familiarity. With a better grasp on which morphemes they saw in training, participants were better able to recognise and reject strings that didn't contain any of these familiar morphemes. As task instructions remained deliberately vague (``do these words belong to what you saw previously?''), participants could interpret the 1M combinations either as incongruent due to the presence of an unknown word part, or congruent due to the presence of a familiar morpheme. In this experiment, participants seems approximately balanced in their preferences for either interpretation.

Women deviated from their slightly more negative evaluation of strings for this 1M condition, possibly due to this ambiguity impacting both genders equally. As for the multilingual measures, the morpheme familiarity effect was strengthened for participants who used their L2 more than their L1 (PC2) and were more balanced across their first top language (L1-L2 difference). Bilinguals were also more sensitive than their monolingual counterparts. This would seem to indicate a kind of bilingual advantage, with language use and balance increasing morphological awareness. This is extended by the result that multilingual experience and multilingual balance (entropy) also lead to a larger morpheme familiarity effect. However, the picture gets blurrier when adding multilingual metrics beyond two languages. Higher L4 scores (PC1) \emph{reduce} morpheme familiarity effects, and quadrilinguals are the multilingual group with the weakest effect. Indeed, looking at the different multilingual categories, quadrilinguals have less variability in ``yes'' responses based on testing condition, followed by the monolinguals, then bilinguals, with trilinguals seeing the strongest effect. This pattern is therefore that of an advantage up to the trilingual level, and then a cost of adding a fourth language. Given the youth of participants (18-25 years age range, 22.21 years on average), having a fourth language in such a short timespan could entail lessened or shallower acquisition of all of a participant's languages due to having to dedicate time to a fourth. This theory would certainly also depend on age of acquisition and use habits of a multilingual's languages. However, since PC1 is sensitive to within-group differences, this analysis also indicates that those quadrilinguals with stronger L4 history/use/proficiency/attitudes have an even larger cost than those with a lesser impact of a fourth language. The pattern of results may also be impacted by the lower statistical power in the quadrilingual group (only 26 quadrilinguals out of 195 participants), but is a peculiar finding that warrants future research.

In the 2M condition, the bilingual metrics of increased L2 compared to L1 use and bilingual balance led to even more sensitivity to the combined presence of two known morphemes, regardless of the combinatorial structure. This reflects the generalised sensitivity discussed above. Multilingual experience also similarly increases ``yes'' responses in the 2M condition. Interestingly, participants with knowledge of languages using more than one script also accept two-morpheme combinations, but this effect is separate from an overarching morpheme familiarity effect. Experience with manipulating linguistic systems with fewer visual differences could enhance visual processing abilities when the stimuli represent a very likely lexical candidate. 

As to whether the stimuli changes freed up enough cognitive resources from the necessary segmentation to allow for more successful co-occurrence tracking, the group-level results do show above-chance performance in the discrimination of within- and between-language morpheme pairings. This is a very small group effect (1\% above chance level), but this was expected by researchers (see experiment 1 \autoref{Participants1} effect size motivations). This is nonetheless significantly above chance in a very novel implicit paradigm, and may represent meaningful, if slight, co-occurrence tracking and categorisation abilities by participants on a group level. This difference was found despite a stronger c bias in the 2M conditions than in previous experimental versions, which is due to the stronger morpheme familiarity effect detailed above and the change in stimulus script. This shows that categorical boundaries can indeed be learned implicitly through repeated exposure to words with co-occurrence regularities. 

Even more interestingly, the accuracy with which participants were able to track these co-occurrences depended on some of the same factors that dictated their sensitivity to the presence of the constituent morphemes. Bilingual participants were the only group to perform above chance level, indicating once more the possibility of a bilingual-specific advantage. Additionally, we find another cost at the quadrilingual level in terms of multilingual balance (cosine similarity). Further information from the balance of multilingual use (use cosine similarity) intimates that this disadvantage may be the consequence of exaggerated \emph{use of a fourth language}. Whereas the PCA in the previous version of the experiment did pick up L4 Use as a valuable component, it was not so in this version, and therefore was not examined in isolation. However, information on the presence and use of a fourth language certainly seems to play a crucial role in mediating the bilingual (and in some case, trilingual) benefit observed here in terms of increased awareness and co-occurrence tracking abilities of morphemes. 


% FAMILIARITY
% better accuracy & faster RTs

% "YES" OVERALL
% gender -> women = less "yes"
% cossim -> more balance = more "yes"

% MORPHEME FAMILIARITY EFFECT
% gender -> women = similar to men only in 1M (otherwise less "yes")
% PC2 (L2 vs L1 use) -> more L2 use = greater morph fam
% L1L2diff -> more bi balance = greater morph fam
% multiexp -> morph multiexp = greater morph fam
% entropy -> higher entropy = greater morph fam
% PC1 (L4) -> lower L4 = greater morph fam
% category -> quadris < monos < bis < tris for morph fam effect strength

% 2M "yes"
% PC2 -> more L2 use = more 2M "yes"
% L1L2diff -> more bi balance = more 2M "yes"
% multiexp -> more multiexp = more 2M "yes"
% script -> mixed script = more 2M "yes"

% 2M ACCURACY
% 2M accuracy above chance - but just by 1%
% higher c than in previous version
% category -> bis = higher 2M scores
% cossim -> higher multi balance = lower 2M scores
% use cossim -> higher multi use balance = lower 2M scores